\FifthArticle{14}
Die vorstehenden Vorschriften lassen sich nun leicht in diejenigen umschmelzen, die der Anwendung der Hilfslogarithmen entsprechen, da sich diese $A=\log a$.
$B=\log b$, $C=\log c$, als die Logarithmen der Quadrate der Tangenten, Kosekanten und Sekanten der von $45^0$ bis $90^0$ zunehmenden Winkel auffassen lassen oder, was dasselbe ist, als die Logarithmen der Quadrate der Kotangenten,
Sekanten und Kosekanten der von $45^0$ bis $0$ abnehmenden Winkel. Also gilt
\[
a=\mathrm{tang}\theta^2,\quad \frac{1}{b}=\sin\theta^2,\quad
\frac{1}{c}=\cos\theta^2
\]
für die Werte von $\theta$ im zweiten Oktanten oder
\[
\frac{1}{a}=\mathrm{tang}\theta^2,\quad
\frac{1}{c}=\sin\theta^2,\quad
\frac{1}{b}=\cos\theta^2
\]
für die Werte von $\theta$ im ersten Oktanten.

Die vollständigen Vorschriften vereinige ich im folgenden Schema, in dem wie oben
\[
\lambda=\frac{f^n}{e^{m+n}}
\]
gesetzt ist.

\Needspace{6\baselineskip}\begin{center}\textit{Erste Form.}\end{center}
\[
x^{m+n}+ex^m-f=0
\]
\textit{Erster Fall.} $\lambda>2^n$
\[
\lambda=a^{m+n}b^n=a^m c^n=\frac{c^{m+n}}{b^m}
\]
\[
x^{m+n}=\frac{f}{b},\quad x^m=\frac{f}{ec},\quad x^n=ea
\]
\textit{Zweiter Fall.} $\lambda<2^n$
\[
\lambda=\frac{b^n}{a^m}=\frac{c^n}{a^{m+n}}=\frac{b^{m+n}}{c^m}
\]
\[
x^{m+n}=\frac{f}{c},\quad x^m=\frac{f}{eb},\quad x^n=\frac{e}{a}
\]

\Needspace{6\baselineskip}\begin{center}\textit{Zweite Form.}\end{center}
\[
x^{m+n}-ex^m-f=0
\]
\textit{Erster Fall.} $\lambda>2^m$
\SourceLeaf{098}
\[
\lambda=a^{m+n}b^m=a^nc^m=\frac{c^{m+n}}{b^n}
\]
\[
x^{m+n}=fb,\quad x^m=\frac{f}{ea},\quad x^n=ec
\]
\textit{Zweiter Fall.}\quad $\lambda<2^m$
\[
\lambda=\frac{b^m}{a^n}=\frac{c^m}{a^{m+n}}=\frac{b^{m+n}}{c^n}
\]
\[
x^{m+n}=fc,\quad x^m=\frac{fa}{e},\quad x^n=eb
\]
\Needspace{6\baselineskip}\begin{center}\textit{Dritte Form.}\end{center}
\[
x^{m+n}-ex^m+f=0
\]
\textit{Erster Fall.}\quad $\frac{1}{\lambda}<\frac{(m+n)^{m+n}}{m^mn^n}$

Gar keine Lösung.

\textit{Zweiter Fall.}\quad $\frac{1}{\lambda}=\frac{(m+n)^{m+n}}{m^mn^n}$

Zwei gleiche Wurzeln, zu deren Bestimmung eine der Gleichungen
\[
x^{m+n}=\frac{fm}{n},\quad x^m=\frac{f(m+n)}{en},\quad x^n=\frac{em}{m+n}
\]
dient.

\textit{Dritter Fall.}\quad $\frac{1}{\lambda}$ größer als $\frac{(m+n)^{m+n}}{m^mn^n}$, aber nicht größer als $2^{m+n}$, und zugleich
$m$ größer als $n$.

Zwei Wurzeln, für die
\[
\frac{1}{\lambda}=a^nb^{m+n}=\frac{c^{m+n}}{a^m}=b^mc^n
\]
\[
x^{m+n}=fa,\quad x^m=\frac{fc}{e},\quad x^n=\frac{e}{b}
\]
\textit{Vierter Fall.}\quad Für $\frac{1}{\lambda}$ dieselben Grenzen wie im dritten Fall, aber $m$ kleiner
als $n$.

Zwei Wurzeln, für die
\[
\frac{1}{\lambda}=a^mb^{m+n}=\frac{c^{m+n}}{a^n}=b^nc^m
\]
\[
x^{m+n}=\frac{f}{a},\quad x^m=\frac{fb}{e},\quad x^n=\frac{e}{c}
\]
\textit{Fünfter Fall.}\quad $\frac{1}{\lambda}$ größer als $2^{m+n}$.

Zwei Wurzeln, von denen die eine durch die Formeln des dritten Falles und
die andere durch die des vierten bestimmt wird.
\SourceLeaf{099}
Es sei noch bemerkt, dass im dritten Fall der Wert von $a$, der der einen Wurzel entspricht, kleiner als $\frac{m}{n}$ und der zur anderen Wurzel gehörende größer als $\frac{m}{n}$ ist; im vierten Fall verhalten sich die beiden Werte von
$a$ in ähnlicher Weise zu $\frac{n}{m}$.

